\documentclass[11pt,letterpaper]{article}
\usepackage{concise-research}
\hypersetup{pdftitle={Shared equipment: a reservation brief},
  pdfauthor={Demonstration},pdfsubject={Synthetic example of a concise research PDF}}
\begin{document}
\brieftitle{Shared equipment: a reservation brief}{Decision and evidence}{Demonstration \quad | \quad Synthetic assumptions}

\textbf{Pilot one shared reservation calendar.} In this fictional workshop,
three duplicate bookings cost 24 minutes each week. A single acceptance
check could prevent those conflicts; the pilot must measure whether people
actually use it. No live system or observed customer data is represented.

\section*{What the example establishes}
The input ledger is deliberately small. It supports an arithmetic estimate,
not a forecast of adoption or financial return. Keep those distinctions in
the final recommendation. Assumptions and calculations appear in
\hyperref[app:inputs]{Appendix A}.

\begin{table}[htbp]
\centering\small
\begin{tabularx}{\linewidth}{@{}lYY@{}}
\toprule
\textbf{Measure} & \textbf{Synthetic input} & \textbf{Decision implication}\\
\midrule
Weekly requests & 60 requests & Small enough for a manual pilot.\\
Duplicate bookings & 3 per week & Define a conflict before counting it.\\
Time per conflict & 8 minutes & 24 minutes of avoidable work per week, if all three are prevented.\\
Adoption & Unknown & Log calendar use; do not assume compliance.\\
\bottomrule
\end{tabularx}
\caption{Illustrative inputs; these are not measured results.}
\end{table}

\section*{A small, explicit workflow}
\begin{figure}[htbp]
\centering
\begin{tikzpicture}[node distance=12mm]
\node[researchbox,text width=27mm] (request) {Request\\resource, time};
\node[researchbox,text width=27mm,right=of request] (check) {Check\\availability};
\node[researchbox,text width=27mm,right=of check] (reserve) {Reserve\\return receipt};
\node[researchbox,text width=27mm,below=of check,dashed] (revise) {Conflict\\choose a slot};
\draw[researcharrow] (request) -- (check);
\draw[researcharrow] (check) -- node[above,font=\footnotesize] {free} (reserve);
\draw[researcharrow,dashed] (check) -- node[right,font=\footnotesize] {occupied} (revise);
\end{tikzpicture}
\caption{Proposed flow. A receipt confirms acceptance; a request alone does not.}
\end{figure}

\section*{How to decide after the pilot}
Record requests, conflicts, and time spent resolving them for one week.
Keep the calendar if conflicts decline without increasing coordination
time. If use is inconsistent, investigate access and ownership before
adding automation. The example makes no causal claim from a one-week comparison.

\clearpage
\appendix
\section{Inputs and a minimal contract}\label{app:inputs}
This appendix demonstrates how to retain implementation detail without
expanding the main brief. Delete it when the reader does not need it.

\subsection*{Reproducible calculation}
All figures are invented for this layout demonstration:
\[
  3\ \mbox{conflicts/week}\times8\ \mbox{minutes/conflict}
  =24\ \mbox{minutes/week}.
\]
The implied conflict rate is $3/60=5\%$. These values do not include
setup time, behavior changes, or partial adoption. They establish an
upper bound on recovered time under the example's assumptions.

\subsection*{Proposed reservation record}
\begin{tabularx}{\linewidth}{@{}lYY@{}}
\toprule
\textbf{Field} & \textbf{Meaning} & \textbf{Constraint}\\
\midrule
\texttt{request\_id} & Caller-generated identifier & Repeated identical requests return the same receipt.\\
\texttt{equipment\_id} & Shared resource & Must identify an available resource.\\
\texttt{start}, \texttt{end} & UTC instants & End follows start; accepted intervals cannot overlap for one resource.\\
\texttt{receipt\_id} & Accepted reservation & Returned only after the reservation is committed.\\
\bottomrule
\end{tabularx}
\par\vspace{4pt}

\subsection*{Failure behavior}
Check availability and accept the reservation atomically. A conflicting
request returns a conflict result without a receipt. Reusing a request ID
with changed content returns an invalid-request result. A timeout leaves
the outcome unknown: look up the original request ID before creating a
new one. These are proposed contracts, not claims about an existing service.

\subsection*{Evidence and links}
\textbf{Source:} the synthetic input ledger in Table 1 and the calculation
above. No external research supports the illustrative numbers.
This \href{https://example.org}{demonstration external link} exists only
to exercise PDF link annotations. Replace it with descriptive links to
actual evidence in a real brief; include the exact section or page when
it matters. The \hyperref[app:inputs]{internal appendix link} demonstrates
cross-reference resolution.
\end{document}
